\documentclass[11pt,letterpaper]{article}

\usepackage[top=1in, bottom=1in, left=1in, right=1in, headheight=14pt]{geometry}
\usepackage{fontspec}
\setmainfont{DejaVu Serif}
\setsansfont{DejaVu Sans}
\setmonofont{DejaVu Sans Mono}

\usepackage{xcolor}
\usepackage{titlesec}
\usepackage{fancyhdr}
\usepackage{enumitem}
\usepackage{hyperref}
\usepackage{booktabs}
\usepackage{tabularx}
\usepackage{longtable}
\usepackage{colortbl}
\usepackage{multirow}
\usepackage{array}
\usepackage{parskip}
\usepackage{tcolorbox}
\usepackage{graphicx}
\usepackage{lastpage}
\usepackage{microtype}

% ── Color Scheme: Music History domain ────────────────────────────────────────
\definecolor{primary}{HTML}{1A237E}       % Deep navy — jazz clubs, night
\definecolor{secondary}{HTML}{880E4F}     % Deep cranberry — Sub Pop red, vinyl
\definecolor{tertiary}{HTML}{2E7D32}      % PNW forest green
\definecolor{accent}{HTML}{1565C0}        % Seattle sky blue
\definecolor{lightprimary}{HTML}{E8EAF6}
\definecolor{lightsecondary}{HTML}{FCE4EC}
\definecolor{lighttertiary}{HTML}{E8F5E9}
\definecolor{rowgray}{HTML}{F5F5F5}
\definecolor{bordergray}{HTML}{BDBDBD}
\definecolor{subtlegray}{HTML}{757575}
\definecolor{darktext}{HTML}{212121}
\definecolor{warngold}{HTML}{F9A825}
\definecolor{blockred}{HTML}{B71C1C}

% ── Typography ─────────────────────────────────────────────────────────────────
\titleformat{\section}
  {\Large\bfseries\sffamily\color{primary}}
  {\thesection}{1em}{}
  [\vspace{-0.5em}\textcolor{bordergray}{\rule{\textwidth}{0.4pt}}]

\titleformat{\subsection}
  {\large\bfseries\sffamily\color{secondary}}
  {\thesubsection}{1em}{}

\titleformat{\subsubsection}
  {\normalsize\bfseries\sffamily\color{tertiary}}
  {\thesubsubsection}{1em}{}

\titlespacing*{\section}{0pt}{1.5em}{0.8em}
\titlespacing*{\subsection}{0pt}{1.2em}{0.5em}
\titlespacing*{\subsubsection}{0pt}{0.8em}{0.3em}

% ── Custom Environments ────────────────────────────────────────────────────────
\tcbuselibrary{breakable,skins}

\newcommand{\stageheader}[2]{%
  \noindent\colorbox{#2}{\makebox[\textwidth][l]{%
    \textcolor{white}{\bfseries\sffamily\large\hspace{8pt}#1\hspace{8pt}}}}%
  \vspace{0.6em}\par%
}

\newtcolorbox{asciibox}{
  colback=rowgray, colframe=bordergray,
  arc=1pt, boxrule=0.5pt,
  left=6pt, right=6pt, top=4pt, bottom=4pt,
  fontupper=\small\ttfamily,
  breakable
}

\newtcolorbox{quotebox}{
  colback=lightprimary, colframe=primary,
  arc=2pt, boxrule=0.8pt,
  left=12pt, right=12pt, top=8pt, bottom=8pt,
  breakable
}

\newtcolorbox{warnbox}{
  colback=lightsecondary, colframe=secondary,
  arc=2pt, boxrule=0.8pt,
  left=12pt, right=12pt, top=8pt, bottom=8pt,
  breakable
}

\newtcolorbox{greenbox}{
  colback=lighttertiary, colframe=tertiary,
  arc=2pt, boxrule=0.8pt,
  left=12pt, right=12pt, top=8pt, bottom=8pt,
  breakable
}

\newcommand{\metafield}[2]{\textbf{\sffamily\textcolor{primary}{#1}} \textcolor{darktext}{#2}\par}

% ── Table helpers ──────────────────────────────────────────────────────────────
\newcolumntype{L}[1]{>{\raggedright\arraybackslash}p{#1}}
\newcolumntype{C}[1]{>{\centering\arraybackslash}p{#1}}
\newcommand{\tblhdr}[1]{\textbf{\sffamily\textcolor{white}{#1}}}

% ── Headers / Footers ─────────────────────────────────────────────────────────
\pagestyle{fancy}
\fancyhf{}
\fancyhead[L]{\small\sffamily\textcolor{subtlegray}{Seattle 360 — Continuous Artist Release Engine}}
\fancyhead[R]{\small\sffamily\textcolor{subtlegray}{GSD Mission Package}}
\fancyfoot[C]{\small\sffamily\textcolor{subtlegray}{Page \thepage\ of \pageref{LastPage}}}
\renewcommand{\headrulewidth}{0.4pt}
\renewcommand{\footrulewidth}{0pt}

% ── Hyperlinks ─────────────────────────────────────────────────────────────────
\hypersetup{
  colorlinks=true,
  linkcolor=accent,
  urlcolor=accent,
  pdfauthor={GSD Ecosystem / Miles Tiberius Foxglove},
  pdftitle={Seattle 360 -- Continuous Artist Release Engine},
  pdfsubject={Vision to Mission Pipeline},
}

% ═══════════════════════════════════════════════════════════════════════════════
\begin{document}

% ── Title Page ─────────────────────────────────────────────────────────────────
\thispagestyle{empty}
\begin{center}
\vspace*{2.5cm}

{\small\sffamily\textcolor{primary}{\textbf{GSD RESEARCH MISSION PACKAGE}}}

\vspace{0.3cm}
\textcolor{bordergray}{\rule{9cm}{0.4pt}}

\vspace{1.2cm}
{\fontsize{30}{36}\selectfont\bfseries\sffamily\textcolor{primary}{Seattle 360}}

\vspace{0.4cm}
{\fontsize{16}{20}\selectfont\bfseries\sffamily\textcolor{secondary}{Continuous Artist Release Engine}}

\vspace{0.6cm}
{\Large\sffamily\textcolor{tertiary}{360 Artists $\cdot$ 360 Degrees $\cdot$ One Unit Circle}}

\vspace{0.4cm}
{\large\sffamily\itshape\textcolor{subtlegray}{A Fleet Research Mission $\cdot$ College of Knowledge Integration}}

\vspace{2.5cm}
{\sffamily\textcolor{accent}{Vision $\rightarrow$ Research $\rightarrow$ Mission}}\\[0.4em]
{\small\sffamily\textcolor{accent}{Full Pipeline $\cdot$ Fleet Activation Profile $\cdot$ Autonomous Continuous Release}}

\vspace{2cm}
\begin{tabular}{ll}
\textcolor{subtlegray}{\sffamily\small Date:} & \textcolor{darktext}{\sffamily\small 2026-03-28} \\[4pt]
\textcolor{subtlegray}{\sffamily\small Input:} & \textcolor{darktext}{\sffamily\small seattle\_360\_geo.csv — 360 artists, degrees 0°–359°} \\[4pt]
\textcolor{subtlegray}{\sffamily\small Activation Profile:} & \textcolor{darktext}{\sffamily\small Fleet (9 components, 2 parallel Wave-1 tracks)} \\[4pt]
\textcolor{subtlegray}{\sffamily\small Model Split:} & \textcolor{darktext}{\sffamily\small $\approx$30\% Opus / $\approx$60\% Sonnet / $\approx$10\% Haiku} \\[4pt]
\textcolor{subtlegray}{\sffamily\small Estimated Budget:} & \textcolor{darktext}{\sffamily\small $\approx$22M tokens across $\approx$720 context windows} \\[4pt]
\textcolor{subtlegray}{\sffamily\small Safety-Critical Tests:} & \textcolor{darktext}{\sffamily\small 22 mandatory-pass gates per release} \\[4pt]
\textcolor{subtlegray}{\sffamily\small Status:} & \textcolor{darktext}{\sffamily\small Ready for GSD Orchestrator} \\
\end{tabular}

\vspace{1.8cm}
\textcolor{bordergray}{\rule{9cm}{0.4pt}}\\[0.4em]
{\small\sffamily\textcolor{subtlegray}{GSD Ecosystem $\cdot$ github.com/Tibsfox $\cdot$ Miles Tiberius Foxglove}}
\end{center}
\newpage

% ── Table of Contents ──────────────────────────────────────────────────────────
\tableofcontents
\newpage

% ═══════════════════════════════════════════════════════════════════════════════
%  STAGE 1 — VISION DOCUMENT
% ═══════════════════════════════════════════════════════════════════════════════
\stageheader{STAGE 1 \quad VISION DOCUMENT}{primary}

\section{Seattle 360 — Vision Guide}

\metafield{Date:}{2026-03-28}
\metafield{Status:}{Initial Vision / Pre-Research}
\metafield{Depends on:}{\texttt{gsd-skill-creator-analysis.md}, \texttt{unit-circle-skill-creator-synthesis.md}, \texttt{.college/} (Rosetta Core), \textit{The Space Between}}
\metafield{Context:}{An autonomous, continuous-release educational engine that processes 360 Seattle-area artists arranged as a unit circle and produces a full Fleet research mission pack per artist, with deep College of Knowledge integration, music theory curriculum mapping, and retrospective-driven learning across all 360 releases.}

\subsection{Vision}

There is a moment every music teacher knows: a student finally \textit{hears} the ii--V--I in a
Bill Evans chord, or feels the tritone substitution drop in a Quincy Jones arrangement,
and something irreversible happens. Theory stops being notation on a page and becomes a
living thing --- audible, tactile, \textit{theirs}. The Seattle music scene, spanning jazz
institutions on the Central District's Jackson Street corridor to the grunge detonation of
1991 Capitol Hill, is one of the richest classrooms on the continent. Every artist in
this dataset is a lesson hiding inside a discography.

The problem is that these lessons have never been taught at scale with the depth they
deserve. Existing resources treat artists as cultural artifacts --- biographical, historical,
anecdotal. They do not trace how Bill Frisell's chromaticism descends from tritone
substitutions Quincy Jones was writing in the Central District forty years earlier. The
connections --- the \textit{spaces between} --- remain invisible.

The Seattle 360 Continuous Release Engine changes that. It treats the 360-artist CSV as
what it structurally is: \textbf{a unit circle}. Degree 0 (Quincy Jones, jazz, low energy, 1950s)
to degree 359 (Unwound, post-hardcore, peak intensity, 2001). Each degree is a release.
Each release is a full Fleet research mission --- sourced, theorized, cross-referenced,
curriculum-mapped --- linking outward to the College of Knowledge's music theory
departments, mathematics foundations from \textit{The Space Between}, and the deep tissue of
Seattle cultural history. Retrospectives after each release accumulate into a living
\texttt{knowledge-state.json} that seeds the next mission, making the 360th release
richer than the first by the accumulated wisdom of 359 prior lessons.

This is the Amiga Principle as pedagogy: not a brute-force encyclopedia of 360 artist
bios, but an architectural system that generates insight through intelligent traversal.

\subsection{Problem Statement}

\begin{enumerate}[leftmargin=2em]
  \item \textbf{Music education is disconnected from actual artists.} Theory curricula teach
    abstract concepts; artist studies teach biography. The bridge --- \textit{show me the theory
    in the music I love} --- is almost never built at scale.

  \item \textbf{Seattle's musical lineage is pedagogically unmapped.} The direct lineage
    from Ray Charles's Central District residency through Jimi Hendrix's blues roots through
    grunge's power-chord minimalism to Bikini Kill's riot grrrl deconstruction is one of
    music history's richest theory progressions. It has never been traced as a curriculum.

  \item \textbf{Research missions are one-and-done.} Each mission is produced in isolation.
    The 200th release has no memory of what the 50th release taught. Lessons decay rather
    than compound.

  \item \textbf{Cross-referencing between College departments is manual.} A music theory node
    in \texttt{.college/music/} and a frequency-ratio proof in \texttt{.college/mathematics/}
    are not automatically linked even when they share deep structure.

  \item \textbf{Cultural attribution is uneven.} Indigenous and Black American musical
    foundations are sources, not subjects, in most curricula. The Central District's jazz
    lineage requires OCAP\textregistered{}-compliant handling and nation-specific attribution.
\end{enumerate}

\subsection{Core Concept}

\textbf{Research $\rightarrow$ Map $\rightarrow$ Link $\rightarrow$ Release $\rightarrow$ Retrospect $\rightarrow$ Carry Forward $\rightarrow$ Repeat.}

The engine reads one artist from the CSV, runs a full Fleet research mission, maps the
artist's musical DNA to College of Knowledge curriculum nodes, creates deep cross-reference
links, produces the release artifacts, runs a retrospective, extracts lessons learned, and
seeds the next artist's mission with accumulated context. The loop repeats 360 times.

\subsubsection{Release Cycle Architecture}

\begin{asciibox}
CSV Row [N]
    |
    v
+--------------------------------------------------+
| WAVE 0: FOUNDATION                               |
| +-----------------+  +------------------------+  |
| | ArtistProfile   |  | KnowledgeState[N]      |  |
| | (from CSV row)  |  | (from retro N-1)       |  |
| +-----------------+  +------------------------+  |
+--------------------------------------------------+
    |                          |
    v                          v
+------------------+   +------------------------+
| WAVE 1A          |   | WAVE 1B                |
| Music Theory     |   | College Knowledge      |
| Mapper (Sonnet)  |   | Linker (Sonnet)        |
| -> TheoryNodes   |   | -> CollegeLinks        |
+-------+----------+   +----------+-------------+
        +---------------+---------+
                        v
        +-------------------------------+
        | WAVE 2: FLEET RESEARCH GEN    |
        | Full research mission PDF     |
        | .tex source + index.html      |
        | (Opus -- synthesis)           |
        +---------------+---------------+
                        v
        +-------------------------------+
        | WAVE 3: RELEASE + RETRO       |
        | Safety Warden (BLOCK gate)    |
        | Publish artifacts             |
        | Retrospective.md              |
        | KnowledgeState[N+1]           |
        +---------------+---------------+
                        v
                  CSV Row [N+1]
\end{asciibox}

\subsubsection{The 360\textdegree{} Unit Circle Mapping}

\begin{asciibox}
                   0 deg: Quincy Jones (Jazz, energy=3)
                          ^
              Jazz        |        Folk/Singer-Songwriter
          (0-22 deg)      |           (45-89 deg)
                          |
  270 deg <---------------+---------------> 90 deg
  Grunge/Metal            |            Indie/Electronic
  (296-359 deg)           |            (86-140 deg)
                          |
                          v
                   180 deg: Blues/Soul/R&B
                   (25-45 deg center mass)

  Energy 1 (innermost) ====> Energy 10 (outermost)
  Radius = intensity of artistic statement
  Arc position = genre/era
  Every degree = one educational release
\end{asciibox}

\subsection{Research Modules}

\subsubsection{Module 1: CSV Intake Engine}
\textbf{Purpose:} Transform each CSV row into a validated \texttt{ArtistProfile}. Derives
\texttt{curriculumDepth} from energy level (1--3 = foundational, 4--6 = intermediate,
7--10 = advanced). Handles CSV anomalies (missing lat/lon, parenthetical name qualifiers).
\textbf{College cross-ref:} None (infrastructure only).

\subsubsection{Module 2: Music Theory Mapper}
\textbf{Purpose:} Map genre + era + energy to a \texttt{TheoryNodeList} --- structured
music theory concepts directly audible in the artist's work, each linked to College
curriculum paths and the running theory genealogy.
\textbf{Key domains:} Jazz harmony, blues form, grunge riff theory, mathcore rhythm, riot
grrrl deconstruction, modal theory, counterpoint, drone/synthesis.
\textbf{College cross-ref:} \texttt{.college/music/theory/}, \texttt{.college/mathematics/}

\subsubsection{Module 3: Fleet Research Generator}
\textbf{Purpose:} Synthesize all Wave 1 outputs into a full three-stage LaTeX research
mission document per artist: Vision Guide + Research Reference + Mission Package.
\textbf{NASA SE verbosity:} Every claim sourced. Every number attributed.
\textbf{College cross-ref:} All departments (writes to Stage 3 integration plan).

\subsubsection{Module 4: College Knowledge Linker}
\textbf{Purpose:} Generate deep links into \texttt{.college/} for each artist --- specifying
exactly which nodes are CREATEd vs. ENRICHed and what content the artist contributes.
The mathematics bridge (unit circle, Fourier, integer sequences) is mandatory for harmonic artists.
\textbf{College cross-ref:} All departments.

\subsubsection{Module 5--8: Release + Safety + Retrospective + Carry-Forward}
\textbf{Purpose:} Publish atomically, enforce Safety Warden gates, extract NASA SE lessons
learned, and write \texttt{KnowledgeState[N+1]} to seed the next mission. These four components
form the Wave 3 critical path. The Safety Warden cannot be bypassed.
\textbf{College cross-ref:} Updates \texttt{college-node-index.json} after each release.

\subsection{Chipset Configuration}

\begin{asciibox}
name: seattle-360-release-engine
version: 1.0.0
description: >
  Autonomous continuous-release engine. 360 artists. One unit circle.
  Fleet activation. Retrospective-driven knowledge accumulation.

agents:
  topology: "pipeline"
  agents:
    - name: "RAVEN"         # Orchestrator -- flight director
    - name: "SALMON"        # CSV intake -- ArtistProfile generation
    - name: "CEDAR"         # Theory mapper -- music curriculum (Wave 1A)
    - name: "ORCA"          # College linker -- deep-link generation (Wave 1B)
    - name: "HAWK"          # Fleet research generator -- Opus synthesis (Wave 2)
    - name: "WARDEN"        # Safety warden -- BLOCK authority (Wave 3a)
    - name: "HERON"         # Release pipeline -- atomic publish (Wave 3b)
    - name: "OWL"           # Retrospective engine -- lessons learned (Wave 3c)
    - name: "DOUGLAS-FIR"   # Carry-forward -- KnowledgeState management (Wave 3d)
    - name: "CAPCOM"        # Human-in-loop -- BLOCK escalation reviews

evaluation:
  gates:
    pre_release:
      - check: "safety-warden-pass"
        action: "block"
      - check: "xelatex-compile-success"
        action: "block"
      - check: "attribution-complete"
        action: "block"
    post_release:
      - check: "knowledge-state-committed"
        action: "block"
      - check: "retrospective-written"
        action: "block"
\end{asciibox}

\subsection{Scope Boundaries}

\subsubsection{In Scope (v1.0)}
\begin{itemize}[leftmargin=2em]
  \item All 360 artists in \texttt{seattle\_360\_geo.csv} (degrees 0--359)
  \item Full Fleet research mission PDF per artist (LaTeX + .tex + index.html)
  \item Music theory curriculum mapping (genre-specific, era-appropriate)
  \item Deep College of Knowledge links (minimum 3 departments per release)
  \item Sequential release with carry-forward retrospectives
  \item Safety Warden enforcement on every release (OCAP\textregistered{}, CARE, UNDRIP)
  \item Progress resumability (interrupt and resume at any artist)
  \item \texttt{knowledge-state.json} accumulation across all 360 releases
\end{itemize}

\subsubsection{Out of Scope (Future)}
\begin{itemize}[leftmargin=2em]
  \item Audio analysis / MIDI representation of theory concepts (v2.0)
  \item Interactive College module rendering in GSD-OS desktop (v2.0)
  \item Seattle $\rightarrow$ Chicago $\rightarrow$ New Orleans lineage comparison (v2.0)
  \item Foxy's Playground Minecraft integration for music-world nodes (v3.0)
\end{itemize}

\subsection{Success Criteria}

\begin{enumerate}[leftmargin=2em]
  \item All 360 artists produce a complete release set in \texttt{releases/seattle-360/NNN-[slug]/}
  \item Every release PDF contains sourced biography, $\geq$3 theory concepts with audible evidence,
        $\geq$3 College of Knowledge cross-references, and a carry-forward section.
  \item The Safety Warden passes every release (no BLOCK-level findings at publication).
  \item \texttt{knowledge-state.json} is updated atomically after every retrospective; engine
        can be interrupted and resumed without state loss.
  \item Every theory concept taught references at least one audible example (track or structural description).
  \item Central District jazz artists carry explicit attribution to Black American musical foundations.
  \item Indigenous cultural connections use nation-specific language (Duwamish, Muckleshoot, Suquamish).
  \item College of Knowledge nodes are tracked in \texttt{knowledge-nodes.json}; no duplicate creation.
  \item By degree 359, \texttt{theory-genealogy.json} traces the complete lineage from jazz to post-hardcore.
  \item Carry-forward demonstrably enriches later missions (genealogy cross-references verifiable).
  \item Every retrospective contains $\geq$1 IMMEDIATE carry item (CONFIRMED, not INFERRED).
  \item Full 360-release run completes within estimated token budget ($\pm$20\%).
\end{enumerate}

\subsection{Relationship to Other Vision Documents}

\begin{center}
\rowcolors{2}{rowgray}{white}
\begin{tabularx}{\textwidth}{L{6cm}X}
\toprule
\rowcolor{primary}
\tblhdr{Document} & \tblhdr{Relationship} \\
\midrule
\texttt{gsd-skill-creator-analysis.md} & Execution engine — skill-creator runs every fleet mission \\
\texttt{unit-circle-skill-creator-synthesis.md} & Mathematical metaphor — 360\textdegree{} CSV IS the unit circle \\
\texttt{gsd-bbs-educational-pack-vision.md} & Output format compatibility \\
\texttt{gsd-chipset-architecture-vision.md} & Chipset YAML patterns for each artist mission \\
\textit{The Space Between} & Mathematical spine — frequency, rhythm, L-systems seeded in theory nodes \\
\texttt{gsd-os-desktop-vision.md} & GSD-OS delivers College of Knowledge content the engine produces \\
\texttt{gsd-foxfire-heritage-skills-vision.md} & Cultural heritage framing for attribution standards \\
\bottomrule
\end{tabularx}
\end{center}

\subsection{The Through-Line}

Jackson Street in Seattle's Central District was once called the ``Harlem of the West'' ---
a corridor where Ray Charles, Quincy Jones, Ernestine Anderson, and dozens more learned
their craft before any of them were famous. That corridor produced a theory of music
that flowed forward through decades: into the ECM jazz of Bill Frisell, into the
soul-inflected grunge that made Sub Pop a household name, into the riot grrrl
deconstructions of Olympia, into the mathcore precision of Botch. The thread never broke.
It just changed key.

This engine is built to trace that thread. The College of Knowledge is the destination;
\textit{The Space Between}'s mathematics are the scaffold; the Amiga Principle is the guide.
Remarkable educational depth not through brute-force biography, but through architectural
intelligence --- each artist a node in a graph where the edges carry the real lessons.
The 360\textdegree{} rotation is complete when the knowledge state has closed the loop: when the
engine can show, in verifiable connected steps, how Quincy Jones at 0\textdegree{} and Unwound at
359\textdegree{} are teaching the same mathematics in different languages.

\newpage

% ═══════════════════════════════════════════════════════════════════════════════
%  STAGE 2 — RESEARCH REFERENCE
% ═══════════════════════════════════════════════════════════════════════════════
\stageheader{STAGE 2 \quad RESEARCH REFERENCE}{secondary}

\section{Research Reference}

\metafield{Purpose:}{Sourced factual foundation for agent execution. Where the vision says
``Central District jazz lineage,'' this document provides verifiable history, specific
artists, dates, and institutions. Where the vision says ``theory mapping,'' this document
provides the genre $\rightarrow$ curriculum mapping tables agents will implement.}

\subsection{How to Use This Document}

\begin{itemize}[leftmargin=2em]
  \item \textbf{Component 02 (Theory Mapper):} Read \S2.3 (Genre Theory Maps) and \S2.4 (Era Vocabulary)
  \item \textbf{Component 03 (Fleet Research Generator):} Read all sections; \S2.2 provides biographical sourcing
  \item \textbf{Component 04 (College Linker):} Read \S2.5 (College Dept Mapping) and \S2.6 (Mathematics)
  \item \textbf{Component 08 (Safety Warden):} Read \S2.7 (Cultural Attribution) and \S2.8 (Sensitivity)
\end{itemize}

\textbf{Key source organizations:}
\begin{itemize}[leftmargin=2em]
  \item \textbf{Seattle Civil Rights \& Labor History Project (UW)} --- Authoritative for Central District jazz history and Black American cultural institutions
  \item \textbf{KEXP (90.3 FM)} --- Seattle's primary music archive; programming records and artist interviews
  \item \textbf{Sub Pop Records} --- Discographic source for grunge, indie, and alt-country releases
  \item \textbf{K Records} --- Source for Olympia-area artists (riot grrrl, lo-fi, indie folk)
  \item \textbf{Grove Music Online / Oxford Music Online} --- Peer-reviewed music theory and biography
  \item \textbf{Black Music Research Journal} --- Scholarly Black music studies (Columbia College Chicago)
\end{itemize}

\subsection{Seattle Music History (Sourced Timeline)}

\begin{center}
\rowcolors{2}{rowgray}{white}
\begin{tabularx}{\textwidth}{L{1.8cm}L{5.2cm}L{3.5cm}X}
\toprule
\rowcolor{primary}
\tblhdr{Era} & \tblhdr{Key Development} & \tblhdr{CSV Degrees} & \tblhdr{Primary Source} \\
\midrule
1930s--40s & Jackson Street jazz corridor & 12 & UW SCRLHP \\
1950s & Garfield H.S. cohort; national labels & 0, 3, 6 & Grove Music Online \\
1960s & Soul/R\&B; Civil Rights context & 27, 30, 34 & Black Music Research Journal \\
1970s & Post-Seattle Act; punk arrives & 31, 32 & KEXP Archive \\
1982--88 & Sub Pop founded; proto-grunge & 294, 297, 310 & Sub Pop Records \\
1988--92 & Grunge explosion; major labels & 311--314 & AllMusic \\
1990--98 & Riot grrrl Olympia; indie folk & 276, 277, 53 & K Records \\
2000--pres. & Post-grunge diversification & 58, 86, 44 & KEXP \\
\bottomrule
\end{tabularx}
\end{center}

\subsection{Genre $\rightarrow$ Music Theory Curriculum Maps}

\subsubsection{Jazz (Degrees 0--22)}

\begin{center}
\rowcolors{2}{rowgray}{white}
\begin{tabularx}{\textwidth}{L{3.8cm}XL{3.2cm}L{2.2cm}}
\toprule
\rowcolor{primary}
\tblhdr{Concept} & \tblhdr{Specificity} & \tblhdr{Example Artist (deg.)} & \tblhdr{Level} \\
\midrule
ii--V--I progression & Diatonic function harmony & Quincy Jones (0) & Intermediate \\
Tritone substitution & Chromatic bII7 for V7 & Bill Frisell (1) & Advanced \\
Modal jazz & Dorian, Mixolydian, Lydian & Marc Seales (13) & Intermediate \\
Chord extensions & Tertian harmony beyond 7th & Thomas Marriott (15) & Intermediate \\
Swing rhythm & Compound triplet subdivision & Floyd Standifer (4) & Foundational \\
Call-and-response & African retention; antiphony & Ernestine Anderson (3) & Foundational \\
Blue notes & Micro-tonal inflection & Overton Berry (5) & Foundational \\
Metric modulation & Tempo within tempo (fusion) & Jeff Lorber (16) & Advanced \\
\bottomrule
\end{tabularx}
\end{center}

\subsubsection{Blues (Degrees 35--39)}

\begin{center}
\rowcolors{2}{rowgray}{white}
\begin{tabularx}{\textwidth}{L{3.8cm}XL{3.2cm}L{2.2cm}}
\toprule
\rowcolor{primary}
\tblhdr{Concept} & \tblhdr{Specificity} & \tblhdr{Example Artist (deg.)} & \tblhdr{Level} \\
\midrule
12-bar blues form & Tonic-subdominant-dominant cycle & Robert Cray (35) & Foundational \\
Blues scale & Pentatonic + $\flat$5 tritone & Ayron Jones (37) & Foundational \\
Bent notes / vibrato & Expressive micro-tonality & Jimi Hendrix (36) & Intermediate \\
Lydian mode & Bright \#4 scale (``Bold as Love'') & Jimi Hendrix (36) & Advanced \\
Dominant 7th as tonic & Non-resolving dominant color & Robert Cray (35) & Intermediate \\
\bottomrule
\end{tabularx}
\end{center}

\subsubsection{Grunge / Metal (Degrees 296--359)}

\begin{center}
\rowcolors{2}{rowgray}{white}
\begin{tabularx}{\textwidth}{L{3.8cm}XL{3.8cm}L{2.2cm}}
\toprule
\rowcolor{primary}
\tblhdr{Concept} & \tblhdr{Specificity} & \tblhdr{Example Artist (deg.)} & \tblhdr{Level} \\
\midrule
Drop-D tuning & Lowered 6th string; power chords & Nirvana (311) & Foundational \\
Lydian mode & Kim Thayil's raised 4th tension & Soundgarden (346) & Advanced \\
Phrygian dominant & $\flat$2 scale; ``Down in a Hole'' & Alice in Chains (347) & Advanced \\
Metric modulation & Tempo-feel shift; mathcore & Botch (344, 352) & Advanced \\
Drone theory & Sustained tones; Helmholtz & Melvins (308, 345) & Advanced \\
Odd time signatures & 7/8, 11/8, prime-number meters & Botch (289, 332) & Advanced \\
Riot grrrl deconstruction & Wrong theory as statement & Bikini Kill (276) & Advanced \\
Power chord & Root + 5th; modal ambiguity & Nirvana (311) & Foundational \\
\bottomrule
\end{tabularx}
\end{center}

\subsection{College Department Mapping}

\begin{center}
\rowcolors{2}{rowgray}{white}
\begin{tabularx}{\textwidth}{L{5.5cm}XL{3.5cm}}
\toprule
\rowcolor{primary}
\tblhdr{.college/ Path} & \tblhdr{Content Type} & \tblhdr{Relevant Degrees} \\
\midrule
\texttt{music/theory/harmony/} & Chord construction, progressions & All \\
\texttt{music/theory/rhythm/} & Meter, subdivision, polyrhythm & All \\
\texttt{music/theory/scales/} & Scales, modes, inflections & All genres \\
\texttt{music/theory/form/} & 12-bar blues, AABA, verse-chorus & Genre-specific \\
\texttt{music/history/jazz/} & Jazz history and evolution & 0--22 \\
\texttt{music/history/seattle/} & Seattle-specific music history & All \\
\texttt{music/history/riot-grrrl/} & Feminist punk history & 276, 277, 280, 357 \\
\texttt{music/history/olympia-krecords/} & K Records ecosystem & 49, 50, 276, 277 \\
\texttt{mathematics/number-theory/ratios/} & Interval ratios (3:2 = 5th) & Any harmonic \\
\texttt{mathematics/trigonometry/unit-circle/} & Pitch cycles = circular functions & Modal artists \\
\texttt{mathematics/calculus/fourier/} & Timbre = sum of sine waves & Electronic \\
\texttt{mathematics/sequences/} & Rhythm as integer sequences & Mathcore, jazz \\
\texttt{mind-body/performance/} & Stage presence, practice psychology & energy $\geq$7 \\
\bottomrule
\end{tabularx}
\end{center}

\subsection{The Unit Circle -- Mathematics Bridge}

\begin{warnbox}
\textbf{\sffamily\textcolor{secondary}{Critical Connection: The CSV IS the Unit Circle}}

The 360-artist dataset maps directly to \textit{The Space Between}'s unit circle chapters.
Artist at degree $N$ occupies radian position $N \times \frac{\pi}{180}$. Energy maps to radius.
Genre maps to quadrant. This is not a metaphor --- it is structural identity.

\medskip
Key bridges that MUST be made explicit in every harmonically-rich release:
\begin{itemize}[leftmargin=1.5em,topsep=2pt]
  \item Perfect 5th = 3:2 frequency ratio $\rightarrow$ \texttt{.college/mathematics/number-theory/ratios/}
  \item Modal theory = scale rotations = circular permutations $\rightarrow$ \texttt{.college/mathematics/trigonometry/unit-circle/}
  \item Timbre = Fourier series (sum of sinusoids) $\rightarrow$ \texttt{.college/mathematics/calculus/fourier/}
  \item Polyrhythm = integer ratio operations (4:3 = perfect fourth, inverted) $\rightarrow$ \texttt{.college/mathematics/sequences/}
  \item Botch's mathcore: 7/8, 11/8 = prime-number meters $\rightarrow$ \texttt{.college/mathematics/sequences/}
\end{itemize}
\end{warnbox}

\subsection{Cultural Attribution Standards}

\begin{center}
\rowcolors{2}{rowgray}{white}
\begin{tabularx}{\textwidth}{L{3.5cm}XL{2.5cm}}
\toprule
\rowcolor{primary}
\tblhdr{Situation} & \tblhdr{Required Language} & \tblhdr{Boundary} \\
\midrule
Describing jazz origins & ``Black American musical tradition rooted in\ldots'' & ABSOLUTE \\
Blues in rock context & ``Blues tradition created by Black American musicians'' & ABSOLUTE \\
Garfield High School & ``Historically Black institution; produced\ldots'' & GATE \\
Indigenous territory & Nation-specific (Duwamish, Muckleshoot, Suquamish) & ABSOLUTE \\
Riot grrrl context & Feminist/political framing explicit & GATE \\
Mia Zapata (The Gits) & Factual, dignified; no sensationalism & ABSOLUTE \\
Living artist claim & Public record only; no speculation & ABSOLUTE \\
\bottomrule
\end{tabularx}
\end{center}

\subsection{Source Bibliography}

\textbf{Authoritative Music Sources:}
\begin{itemize}[leftmargin=2em,topsep=2pt]
  \item Grove Music Online / Oxford Music Online --- peer-reviewed music encyclopedia
  \item AllMusic database --- discographic data, genre classifications
  \item Black Music Research Journal (Columbia College Chicago)
  \item Journal of the Society for American Music (peer-reviewed)
\end{itemize}

\textbf{Seattle-Specific:}
\begin{itemize}[leftmargin=2em,topsep=2pt]
  \item Seattle Civil Rights \& Labor History Project (UW): \url{depts.washington.edu/civilr/}
  \item KEXP programming archives: \url{kexp.org}
  \item Sub Pop Records discography: \url{subpop.com}
  \item K Records catalog: \url{krecs.com}
\end{itemize}

\textbf{Cultural Attribution:}
\begin{itemize}[leftmargin=2em,topsep=2pt]
  \item OCAP\textregistered{} Principles (First Nations Information Governance Centre): \url{fnigc.ca}
  \item CARE Principles for Indigenous Data Governance: \url{gida-global.org}
  \item UN Declaration on the Rights of Indigenous Peoples (UNDRIP): \url{un.org}
  \item Duwamish Tribe: \url{duwamishtribe.org}
\end{itemize}

\newpage

% ═══════════════════════════════════════════════════════════════════════════════
%  STAGE 3 — MISSION PACKAGE
% ═══════════════════════════════════════════════════════════════════════════════
\stageheader{STAGE 3 \quad MISSION PACKAGE}{tertiary}

\section{Milestone Specification}

\metafield{Mission Objective:}{Build and execute a fully autonomous, continuous-release educational
engine that processes all 360 artists in \texttt{seattle\_360\_geo.csv} sequentially. ``Done'' looks
like \texttt{releases/seattle-360/359-unwound/} existing as a complete, Safety-Warden-passed release,
with \texttt{theory-genealogy.json} tracing the connected theoretical lineage from degree 0
(Quincy Jones) through degree 359 (Unwound).}

\subsection{Architecture Overview}

\begin{asciibox}
INPUT: seattle_360_geo.csv (360 rows, degrees 0-359)
                |
                v
+----------------------------------------------------------+
| WAVE 0: FOUNDATION (Sequential, <5 min per artist)       |
|   CSV Row[N] --> ArtistProfile[N]                        |
|   knowledge-state.json --> ActiveContext[N] (<=2K)       |
|   college-node-index.json --> CollegeIndex[N]            |
|   theory-genealogy.json --> TheoryGenealogy[N]           |
+------------------+-------------------------------+-------+
                   |                               |
                   v                               v
+--------------------+              +---------------------+
| WAVE 1A: Theory    |              | WAVE 1B: College    |
| Mapping (CEDAR)    |              | Linking (ORCA)      |
| Sonnet             |              | Sonnet              |
| -> TheoryNodeList  |              | -> CollegeLinkList  |
+----------+---------+              +-----------+---------+
           +-------------------+---+
                               v
           +-----------------------------------+
           | WAVE 2: FLEET RESEARCH (HAWK/Opus)|
           |   PDF + .tex + index.html + json  |
           +------------------+----------------+
                              v
           +-----------------------------------+
           | WAVE 3a: SAFETY WARDEN (Opus)     |
           |   BLOCK --> CAPCOM --> halt        |
           |   PASS  --> continue              |
           +------------------+----------------+
                              v
           +-----------------------------------+
           | WAVE 3b: RELEASE (HERON/Sonnet)   |
           |   Atomic publish + write-protect  |
           +------------------+----------------+
                              v
           +-----------------------------------+
           | WAVE 3c: RETROSPECTIVE (OWL/Opus) |
           |   NASA SE lessons learned         |
           +------------------+----------------+
                              v
           +-----------------------------------+
           | WAVE 3d: CARRY FORWARD (D-FIR)    |
           |   KnowledgeState[N+1]             |
           |   N<359: loop  |  N=359: COMPLETE |
           +------------------+----------------+
\end{asciibox}

\subsection{Component Breakdown}

\begin{center}
\rowcolors{2}{rowgray}{white}
\begin{tabularx}{\textwidth}{C{0.3cm}L{3.8cm}C{1cm}C{1.2cm}C{1.5cm}C{2cm}L{2.5cm}}
\toprule
\rowcolor{primary}
\tblhdr{\#} & \tblhdr{Component} & \tblhdr{Wave} & \tblhdr{Track} & \tblhdr{Model} & \tblhdr{Tokens/Artist} & \tblhdr{Depends On} \\
\midrule
0 & Shared Types & 0 & --- & Haiku & $\sim$4K (once) & None \\
1 & CSV Intake Engine & 0 & --- & Haiku & $\sim$2K & \#0 \\
2 & Music Theory Mapper & 1A & A & Sonnet & $\sim$8K & \#1 \\
3 & Fleet Research Generator & 2 & --- & Opus & $\sim$20K & \#2, \#4 \\
4 & College Knowledge Linker & 1B & B & Sonnet & $\sim$6K & \#1 \\
5 & Release Pipeline & 3b & --- & Sonnet & $\sim$3K & \#3, \#8 \\
6 & Retrospective Engine & 3c & --- & Opus & $\sim$8K & \#5 \\
7 & Carry-Forward Controller & 3d & --- & Sonnet & $\sim$4K & \#6 \\
8 & Safety Warden & 3a & --- & Opus & $\sim$6.5K & \#3 \\
\midrule
\rowcolor{lightprimary}
& \textbf{Per-artist total} & & & & $\sim$\textbf{57.5K} & \\
& \textbf{360-run total} & & & & $\sim$\textbf{20.7M} & \\
\bottomrule
\end{tabularx}
\end{center}

\subsection{Deliverables Table}

\begin{center}
\rowcolors{2}{rowgray}{white}
\begin{tabularx}{\textwidth}{C{0.3cm}L{4cm}XL{4cm}}
\toprule
\rowcolor{primary}
\tblhdr{\#} & \tblhdr{Deliverable} & \tblhdr{Acceptance Criteria} & \tblhdr{Component Spec} \\
\midrule
1 & \texttt{ArtistProfile[360]} & All 360 parsed; schema-valid & \texttt{components/01-csv-intake-engine.md} \\
2 & \texttt{TheoryNodeList} per artist & $\geq$3 theory concepts; audible evidence & \texttt{components/02-music-theory-mapper.md} \\
3 & Fleet research PDF per artist & XeLaTeX compiles; 3 stages; Warden PASS & \texttt{components/03-fleet-research-generator.md} \\
4 & \texttt{CollegeLinkList} per artist & $\geq$3 links; no duplicates & \texttt{components/04-college-knowledge-linker.md} \\
5 & Release artifact set (360×) & All files in \texttt{NNN-[slug]/} paths & \texttt{components/05-release-pipeline.md} \\
6 & \texttt{retrospective.md} per artist & $\geq$1 IMMEDIATE carry item; sourced & \texttt{components/06-retrospective-engine.md} \\
7 & \texttt{knowledge-state.json} (final) & Atomic; full genealogy; 360 entries & \texttt{components/07-carry-forward-controller.md} \\
8 & Safety audit log (360×) & All releases PASS; BLOCKs escalated & \texttt{components/08-safety-warden.md} \\
\bottomrule
\end{tabularx}
\end{center}

\section{Wave Execution Plan}

\subsection{Wave Summary}

\begin{center}
\rowcolors{2}{rowgray}{white}
\begin{tabularx}{\textwidth}{L{1.5cm}L{4cm}L{2cm}L{2.5cm}C{2cm}X}
\toprule
\rowcolor{primary}
\tblhdr{Wave} & \tblhdr{Name} & \tblhdr{Mode} & \tblhdr{Models} & \tblhdr{Est. Tokens} & \tblhdr{Cache Benefit} \\
\midrule
0 & Foundation & Sequential & Haiku & $\sim$6K & State summary cached \\
1A & Theory Mapping & Parallel & Sonnet & $\sim$8K & ArtistProfile cached \\
1B & College Linking & Parallel & Sonnet & $\sim$6K & ArtistProfile cached \\
2 & Fleet Research & Sequential & Opus & $\sim$20K & Wave 1 outputs cached \\
3a & Safety Audit & Sequential & Opus & $\sim$6.5K & PDF cached \\
3b & Release & Sequential & Sonnet & $\sim$3K & Safety pass cached \\
3c & Retrospective & Sequential & Opus & $\sim$8K & Release record cached \\
3d & Carry Forward & Sequential & Sonnet & $\sim$4K & Retro record cached \\
\bottomrule
\end{tabularx}
\end{center}

\subsection{Dependency Graph}

\begin{asciibox}
[Wave 0: Foundation]
   |- ArtistProfile[N] -------------------------+
   |- ActiveContext[N] (<=2K) ----------------+-|
   |- CollegeIndex[N] ----------------------+-|-|
   |- TheoryGenealogy[N] ------------------+-|-|-|
                                           | | | |
                         +-----------+    | | | |
                         | Track 1A  |<---+ | | |
                         | Theory    |<-----+ | |
                         | Mapper    |        | |
                         +-----+-----+        | |
                               |         +----+-+-+
                               |         | Track 1B|
                               |         | College |
                               |         | Linker  |
                               |         +----+----+
                               +----+----+    |
                                    |         |
                                    v         |
                              [Wave 2: Fleet Research]
                                    |
                             +------+-------+
                             | Wave 3a:     |
                             | Safety Warden|
                             | BLOCK/PASS   |
                             +------+-------+
                          BLOCK |   | PASS
                                |   v
                      CAPCOM -> [Wave 3b: Release]
                                    |
                                    v
                             [Wave 3c: Retro]
                                    |
                                    v
                             [Wave 3d: Carry]
                                    |
                          N<359: loop to Wave 0[N+1]
                          N=359: CYCLE-COMPLETE
\end{asciibox}

\section{Test Plan Summary}

\begin{center}
\rowcolors{2}{rowgray}{white}
\begin{tabularx}{\textwidth}{L{5cm}C{1.5cm}L{2.5cm}X}
\toprule
\rowcolor{primary}
\tblhdr{Category} & \tblhdr{Count} & \tblhdr{Priority} & \tblhdr{Failure Action} \\
\midrule
Safety-critical (SC-01 to SC-22) & 22 & Mandatory & BLOCK release \\
Core functionality (per-component) & 64 & Required & BLOCK release \\
Integration (cross-component) & 32 & Required & BLOCK release \\
Edge cases (unusual artists/genres) & 18 & Best-effort & LOG + continue \\
Performance (token/time budget) & 12 & Best-effort & WARN + log \\
\midrule
\rowcolor{lightprimary}
\textbf{Total} & \textbf{148} & & \\
\bottomrule
\end{tabularx}
\end{center}

\subsubsection{Selected Safety-Critical Tests (BLOCK on Failure)}

\begin{center}
\rowcolors{2}{rowgray}{white}
\begin{tabularx}{\textwidth}{L{1.5cm}XL{3.5cm}}
\toprule
\rowcolor{blockred}
\tblhdr{Test ID} & \tblhdr{Verifies} & \tblhdr{Component} \\
\midrule
SC-01 & All biographical claims have professional citation & 03-fleet-research-gen \\
SC-02 & Indigenous attribution uses nation-specific language & 08-safety-warden \\
SC-03 & Black American musical attribution explicit & 08-safety-warden \\
SC-06 & Mia Zapata dignity protocol (The Gits) & 08-safety-warden \\
SC-09 & XeLaTeX compiles without errors & 03-fleet-research-gen \\
SC-10 & Safety Warden BLOCK prevents publication & 05-release-pipeline \\
SC-11 & Atomic release writes (no partial artifacts) & 05-release-pipeline \\
SC-13 & KnowledgeState atomic write & 07-carry-forward \\
SC-17 & Theory claims have audible evidence & 02-theory-mapper \\
SC-20 & Engine resumability verified & 05-release-pipeline \\
SC-21 & Opus used for Safety Warden (not Sonnet/Haiku) & 08-safety-warden \\
SC-22 & OCAP\textregistered{} compliance for Indigenous knowledge & 08-safety-warden \\
\bottomrule
\end{tabularx}
\end{center}

\section{Mission Crew Manifest}

\begin{center}
\rowcolors{2}{rowgray}{white}
\begin{tabularx}{\textwidth}{L{3cm}L{2.5cm}L{2.2cm}X}
\toprule
\rowcolor{primary}
\tblhdr{Role} & \tblhdr{Agent (PNW)} & \tblhdr{Wave} & \tblhdr{Responsibility} \\
\midrule
FLIGHT & RAVEN & All waves & Flight director; go/no-go authority; wave sequencing \\
PLAN & SALMON & Wave 0 & CSV ingestion; ArtistProfile generation \\
INTEL-A & CEDAR & Wave 1A & Music theory curriculum mapping (Sonnet) \\
INTEL-B & ORCA & Wave 1B & College Knowledge deep-link generation (Sonnet) \\
EXEC & HAWK & Wave 2 & Fleet research mission synthesis (Opus) \\
SAFETY & WARDEN & Wave 3a & Attribution, accuracy, cultural sensitivity (Opus; BLOCK authority) \\
RELEASE & HERON & Wave 3b & Sequential atomic release coordination (Sonnet) \\
RETRO & OWL & Wave 3c & Retrospective lessons learned extraction (Opus) \\
STATE & DOUGLAS-FIR & Wave 3d & KnowledgeState management; carry-forward (Sonnet) \\
CAPCOM & Human & BLOCK gate & Reviews BLOCK-level safety findings; authorizes resume \\
\bottomrule
\end{tabularx}
\end{center}

\section{Execution Summary}

\begin{center}
\rowcolors{2}{rowgray}{white}
\begin{tabularx}{\textwidth}{L{5.5cm}X}
\toprule
\rowcolor{primary}
\tblhdr{Metric} & \tblhdr{Value} \\
\midrule
Total artists & 360 (degrees 0°--359°) \\
Total components & 9 (including Safety Warden) \\
Parallel tracks & 2 (Wave 1A Theory + Wave 1B College) \\
Sequential depth & 4 waves per artist cycle \\
Activation profile & Fleet (Full crew, 10 roles including CAPCOM) \\
Model split & $\approx$30\% Opus / $\approx$60\% Sonnet / $\approx$10\% Haiku \\
Tokens per artist & $\approx$57.5K \\
Total token budget & $\approx$20.7M \\
Estimated context windows & $\approx$720 (2 per artist) \\
Estimated wall time & $\approx$180 sessions (2 artists/session) \\
Total tests & 148 (22 safety-critical mandatory-pass) \\
College departments linked & 7+ per release \\
Output root & \texttt{releases/seattle-360/} \\
State file & \texttt{releases/seattle-360/knowledge-state.json} \\
Genealogy file & \texttt{releases/seattle-360/theory-genealogy.json} \\
Completion signal & \texttt{CYCLE-COMPLETE} after degree 359 (Unwound) \\
\bottomrule
\end{tabularx}
\end{center}

\vspace{2em}
\begin{quotebox}
\textit{``Jackson Street was once the Harlem of the West. The engine completes its
rotation when the knowledge state can show, in verifiable connected steps, how Quincy Jones
at 0° and Unwound at 359° are teaching the same mathematics in different languages.
The thread never broke. It just changed key.''}

\medskip
\textcolor{subtlegray}{\sffamily\small Through-line: Seattle 360 Continuous Release Engine $\cdot$ GSD Ecosystem $\cdot$ 2026}
\end{quotebox}

\end{document}
